\chapter{Large N Conformal Data and Holographic Fields}
\label{ch:volume-iii-large-n-conformal-data-holographic-fields}

\section{Large N Factorization in a Conformal Theory}

Consider a sequence of conformal field theories indexed by an integer \(N\).
Assume the stress-tensor two-point coefficient satisfies
\[
  C_T(N)\longrightarrow\infty
\]
as \(N\to\infty\).  In adjoint gauge theories \(C_T=O(N^2)\).  Choose
single-trace primary operators \(\mathcal O_i\) normalized so that
\[
  \left\langle\mathcal O_i(x)\mathcal O_j(0)\right\rangle
  =
  \frac{\delta_{ij}}{|x|^{2\Delta_i}}
\]
at separated points.  Large \(N\) factorization means that connected
correlators scale as
\begin{equation}
  \left\langle
    \mathcal O_{i_1}\cdots\mathcal O_{i_k}
  \right\rangle_{\rm conn}
  =
  O(C_T^{1-k/2})
  \label{eq:large-n-ct-factorization}
\end{equation}
for fixed \(k\) and fixed operator labels.  With \(C_T\sim N^2\), this is the
usual \(N^{2-k}\) scaling.

Multi-trace operators are composites built from products of single traces and
their derivatives.  At \(N=\infty\), their dimensions are sums
\[
  \Delta_{\mathcal O_i\mathcal O_j,n,\ell}
  =
  \Delta_i+\Delta_j+2n+\ell,
  \qquad
  n\in\mathbb Z_{\ge0},
\]
up to degeneracies and mixing.  At large but finite \(N\), anomalous
dimensions and OPE coefficients of these multi-trace operators are determined
by connected correlators and are expanded in powers of \(1/C_T\).

\begin{definition}[Holographic large \(N\) CFT]
  A holographic large \(N\) conformal field theory, in the sense used here, is
  a large \(C_T\) family with factorization, a parametrically sparse spectrum
  of low-dimension single-trace primaries above the stress tensor and
  conserved currents, and connected correlators that admit an expansion
  organized by local effective interactions in an emergent
  \(\mathrm{AdS}_{D+1}\) spacetime.
\end{definition}

The definition is operational.  It identifies the CFT conditions under which
correlators can be reorganized as amplitudes of weakly coupled fields in
anti-de Sitter space.

\section{Euclidean Anti-de Sitter Space}

Euclidean \(\mathrm{AdS}_{D+1}\), also called hyperbolic space
\(\mathbb H^{D+1}\), is the upper half-space
\[
  z>0,\qquad x\in\mathbb R^D
\]
with metric
\begin{equation}
  \dd s^2
  =
  R^2\frac{\dd z^2+\dd x^i\dd x^i}{z^2},
  \qquad i=1,\ldots,D.
  \label{eq:eads-upper-half-space-metric}
\end{equation}
The conformal boundary is at \(z=0\).  The isometry group of
\eqref{eq:eads-upper-half-space-metric} is \(SO(D+1,1)\), the Euclidean
conformal group acting on the boundary coordinates \(x^i\).

A scalar bulk field \(\varphi\) with action
\begin{equation}
  S_{\rm bulk}[\varphi]
  =
  \int_{\mathrm{AdS}_{D+1}}\dd^Dx\,\dd z\,\sqrt g
  \left[
    \frac12g^{MN}\partial_M\varphi\partial_N\varphi
    +\frac12m^2\varphi^2
  \right]
  +\cdots
  \label{eq:eads-scalar-action}
\end{equation}
satisfies
\[
  \left(-\nabla^2+m^2\right)\varphi=0
\]
at the linearized level.  Near the boundary, solutions behave as
\[
  \varphi(x,z)
  =
  z^{D-\Delta}\varphi_0(x)
  +
  z^\Delta A(x)
  +\cdots ,
\]
where
\begin{equation}
  m^2R^2=\Delta(\Delta-D).
  \label{eq:scalar-mass-dimension-relation}
\end{equation}
The root used for the standard quantization is
\[
  \Delta=\frac D2+\sqrt{\frac{D^2}{4}+m^2R^2}.
\]
When the mass lies in the alternative-quantization window, the other root can
also define a unitary boundary condition; the choice of boundary condition is
part of the theory.

\begin{figure}[H]
  \centering
  \resizebox{0.94\textwidth}{!}{%
  \begin{tikzpicture}[x=1cm,y=1cm,every node/.style={font=\scriptsize}]
    \tikzset{
      arrow/.style={-{Latex[length=2mm]},line width=0.85pt},
      ray/.style={blue!65!black,line width=0.9pt},
      boundary/.style={violet!80!black,line width=1.1pt},
      bulk/.style={fill=blue!4,draw=blue!45!black,line width=0.7pt}
    }
    \fill[blue!3] (-4.8,-1.30) rectangle (4.8,1.45);
    \draw[boundary] (-4.8,1.45)--(4.8,1.45);
    \node[violet!80!black] at (3.95,1.80) {boundary \(z=0\)};
    \draw[->,thick] (-4.45,1.20)--(-4.45,-1.05) node[below] {\(z\)};
    \foreach \x/\lab in {-3.2/{x_1},-1.1/{x_2},1.65/{x_3}} {
      \fill[violet!80!black] (\x,1.45) circle (2.0pt);
      \node[violet!80!black,above] at (\x,1.45) {\(\lab\)};
    }
    \node[bulk,circle,inner sep=2.2pt] (v) at (-0.65,-0.18) {};
    \node[blue!65!black,right] at (v) {\(\varphi(x,z)\)};
    \foreach \x in {-3.2,-1.1,1.65} {
      \draw[ray] (\x,1.45) .. controls (\x,0.55) and (-0.65,0.35) .. (v);
    }
    \node[align=center,green!45!black] at (2.70,-0.80)
      {bulk field is dual to\\boundary primary \(\mathcal O\)};
    \draw[arrow,green!45!black] (1.70,-0.62)--(0.03,-0.22);
  \end{tikzpicture}
  }
  \caption{The boundary value \(\varphi_0(x)\) is the source for a CFT primary
  \(\mathcal O(x)\), while the bulk equation fixes the dependence on the
  radial coordinate \(z\).}
  \label{fig:holographic-boundary-source-bulk-field}
\end{figure}

\section{Generating Functional}

The holographic relation between a bulk field and a scalar primary
\(\mathcal O\) is expressed by the generating functional
\begin{equation}
  Z_{\rm bulk}[\varphi_0]
  =
  \left\langle
    \exp\left(\int\dd^Dx\,\varphi_0(x)\mathcal O(x)\right)
  \right\rangle_{\rm CFT}.
  \label{eq:holographic-generating-functional}
\end{equation}
At large \(C_T\), the left side is evaluated by a saddle:
\[
  Z_{\rm bulk}[\varphi_0]
  =
  \exp\left[-S_{\rm bulk}^{\rm ren}[\varphi_{\rm cl}]\right]
  \left(1+O(C_T^{-1})\right).
\]
Here \(\varphi_{\rm cl}\) solves the bulk equations with boundary condition
\[
  \lim_{z\to0}z^{\Delta-D}\varphi_{\rm cl}(x,z)=\varphi_0(x),
\]
and \(S_{\rm bulk}^{\rm ren}\) is the renormalized on-shell action including
local boundary counterterms.

For the metric \eqref{eq:eads-upper-half-space-metric}, the scalar
boundary-to-bulk propagator is
\begin{equation}
  K_\Delta(x,z;x_0)
  =
  C_\Delta
  \left(
    \frac{z}{z^2+|x-x_0|^2}
  \right)^\Delta,
  \qquad
  C_\Delta=
  \frac{\Gamma(\Delta)}
       {\pi^{D/2}\Gamma(\Delta-D/2)}.
  \label{eq:boundary-to-bulk-propagator}
\end{equation}
The classical solution linear in the source is
\[
  \varphi_{\rm cl}(x,z)
  =
  \int\dd^Dx_0\,K_\Delta(x,z;x_0)\varphi_0(x_0).
\]
Substituting this solution into the quadratic action gives the two-point
function
\begin{equation}
  \left\langle\mathcal O(x_1)\mathcal O(x_2)\right\rangle
  =
  \frac{(2\Delta-D)C_\Delta}{|x_{12}|^{2\Delta}}
  \label{eq:holographic-scalar-two-point-function}
\end{equation}
up to the overall normalization chosen for the bulk kinetic term.  Choosing a
canonical CFT normalization fixes that kinetic coefficient.

\section{Spin, Gauge Fields, and the Stress Tensor}

A conserved spin-one current \(J_i\) in the CFT has dimension
\[
  \Delta_J=D-1.
\]
It is dual to a bulk gauge field \(A_M\).  Boundary gauge transformations of
the source \(A_i^{(0)}\) encode current conservation:
\[
  \nabla^iJ_i=0.
\]
The bulk gauge symmetry is therefore the redundancy of the source description
for a conserved CFT current.

The stress tensor \(T_{ij}\) has dimension
\[
  \Delta_T=D
\]
and is dual to the bulk metric \(g_{MN}\).  The boundary value of the metric
is the CFT background metric up to Weyl rescaling:
\[
  g_{ij}(x,z)
  \sim
  \frac{R^2}{z^2}g^{(0)}_{ij}(x)
  \qquad
  (z\to0).
\]
The CFT generating functional in a curved background is obtained by varying
the renormalized gravitational action with respect to \(g^{(0)}_{ij}\):
\[
  \left\langle T_{ij}(x)\right\rangle
  =
  \frac{2}{\sqrt{g^{(0)}}}
  \frac{\delta \log Z_{\rm bulk}}{\delta g^{(0)\,ij}(x)}.
\]
The coefficient \(C_T\) fixes the effective Newton constant through
\[
  C_T\propto \frac{R^{D-1}}{G_N},
\]
with a dimension-dependent proportionality constant set by the normalization
of \(T_{ij}\).

\begin{figure}[H]
  \centering
  \resizebox{0.88\textwidth}{!}{%
  \begin{tikzpicture}[x=1cm,y=1cm,every node/.style={font=\scriptsize}]
    \tikzset{
      arrow/.style={-{Latex[length=2mm]},line width=0.85pt},
      box/.style={draw,rounded corners=4pt,line width=0.8pt,fill=blue!4,
        minimum width=3.2cm,minimum height=0.86cm,align=center}
    }
    \node[box] (scalar) at (-3.9,1.20)
      {scalar field \(\varphi\)\\\(\leftrightarrow\) primary \(\mathcal O\)};
    \node[box] (gauge) at (-3.9,0)
      {gauge field \(A_M\)\\\(\leftrightarrow\) current \(J_i\)};
    \node[box] (metric) at (-3.9,-1.20)
      {metric \(g_{MN}\)\\\(\leftrightarrow\) stress tensor \(T_{ij}\)};
    \node[box] (mass) at (1.05,1.20)
      {\(m^2R^2=\Delta(\Delta-D)\)};
    \node[box] (sym) at (1.05,0)
      {gauge redundancy\\\(\Rightarrow\) Ward identity};
    \node[box] (gravity) at (1.05,-1.20)
      {\(R^{D-1}/G_N\)\\sets \(C_T\)};
    \draw[arrow] (scalar)--(mass);
    \draw[arrow] (gauge)--(sym);
    \draw[arrow] (metric)--(gravity);
  \end{tikzpicture}
  }
  \caption{The holographic dictionary maps CFT representation-theoretic data
  to bulk fields and their gauge redundancies.}
  \label{fig:holographic-field-dictionary}
\end{figure}

\section{Effective Bulk Expansion}

The bulk description is an effective theory organized by two parameters.  The
first is the loop-counting parameter
\[
  G_N/R^{D-1}\sim C_T^{-1}.
\]
Tree-level Witten diagrams compute the leading connected correlators, and
bulk loops compute corrections suppressed by powers of \(1/C_T\).  The second
is the derivative expansion.  If the single-trace spectrum has a gap
\(\Delta_{\rm gap}\) above the light fields, then local higher-derivative
interactions are suppressed by powers of \(\Delta_{\rm gap}^{-1}\).

For scalar fields \(\varphi^I\), a local bulk Lagrangian has the schematic
form
\[
  \mathcal L_{\rm bulk}
  =
  \frac1{16\pi G_N}(R_{\rm bulk}-2\Lambda)
  +
  \frac12(\nabla\varphi^I)^2
  +
  \frac12m_I^2(\varphi^I)^2
  +
  g_{IJK}\varphi^I\varphi^J\varphi^K
  +
  \sum_{n\ge0}c_n\nabla^{2n}\varphi^4+\cdots .
\]
Each coefficient is another way of encoding CFT data.  The cubic couplings
are fixed by CFT three-point coefficients after normalizations are chosen.
Quartic and higher interactions are constrained by crossing of four-point and
higher-point functions.

The emergent radial direction is the geometric representation of scale in
the CFT.  Boundary points near each other probe the asymptotic region
\(z\to0\); long-distance operator products probe the propagation of bulk
fields through the interior.  This is a precise reorganization of the CFT
operator algebra whenever the large \(N\), sparse-spectrum assumptions are
satisfied.
